\section{Overview}

Generalist ``\emph{foundation}'' robot models are the next step toward the widespread adoption of intelligent and robust robotic systems. 
Inspired by the tremendous advancements in large language, vision, and vision-language models, the development of robotic foundation models is just beginning. 
However, ensuring their security, safety, and trustworthiness remains a significant challenge. Building on these models, system designers need comprehensive frameworks, advanced security measures, and specialized methods to build secure, safe, and reliable models for downstream tasks. This proposal aims to address these critical needs, improving the security and trustworthiness of foundation robot models for real-world deployment. 

Millions of industrial robots (\eg logistics and manufacturing) and service robots (\eg medicine, home, transportation, defense, and military) are poised to become smarter and more reliable with advances in these foundation models. The previous barrier to the emergence of large robotic foundation models, namely, the need for large-scale ``\emph{Internet-wide}'' datasets as seen in large language models (LLMs) and vision-language models (VLMs), has diminished. Recent efforts employ large-scale \xot robotic datasets and new training methods, making this challenge less critical and paving the way for significant advancements.
Current efforts to build robotic foundation models, \eg Language-Vision-Action (VLA) models, focus on integrating various pre-trained models (\eg vision models, LLMs, and VLMs) for multiple modalities (\eg visual and textual instructions) used by robots. 
The transfer learning capabilities of these components enable robotic systems to understand, interact, and learn new tasks. However, this raises concerns about potential security and safety issues that might be inherited from these models, along with other adversarial threats.% they may face.

\vspace{-2ex}
%This proposal seeks to fill the gap in hardware platforms, runtime systems, and tools for practical, dynamic, and capable applications on intermittently powered embedded systems.
\observationdef{
\begin{observation*}{}{}
\vspace{-1ex}\textbf{This proposal seeks to address the security and trustworthiness challenges in
developing foundation models for robotics, ensuring robust frameworks and defense mechanisms for
integrating various pre-trained models, and securing against inherited vulnerabilities and adversarial
threats.}
\vspace{-1ex}
\end{observation*}
}
\vspace{-3ex}

% Neural network-based models are susceptible to adversarial manipulations, where imperceptible perturbations to the input can drastically alter the output (\eg targeting a specific output or a random incorrect output). 
Studying vulnerabilities in vision and LLMs, which have been the theme of AI security research for many researchers including the PIs, does not transfer cleanly to robotics. %Robotics is just beginning to tackle the complexities of securing large models. 
The integration of multiple modalities/model-components to generate physical interaction with the robot environment introduces complex and largely unexplored attack surface that can be exploited to cause harmful behaviors, safety hazards, and operational failures.
Ensuring the safe deployment of robotic technologies across various applications requires a focused effort to understand and address these security challenges.  
A robot policy returns a joint-angle or an end-effector pose, so every manipulated output becomes a
physical displacement rather than a string of textual tokens. Security research on robot models has yet to
account for this distinction, leaving the field without shared threat models, defenses, or benchmarks
grounded in the kinematics of the robot. 
% and without a benchmark comparing results across models and embodiments.
This project closes those gaps by characterizing the threat space of generalist robot models across
embodiments and levels of adversarial knowledge, exposing how vulnerabilities propagate across coupled
modalities, grounding defenses in the action space of the robot, and releasing benchmarks for
safety, robustness, and interpretability. %This makes the physical consequences of adversarial manipulation become measurable, defensible, and reproducible for generalist robot models.
%This work enables the development and adoption of secure and trustworthy robotic models in real-world applications.





\begin{figure*}[t]
  \centering
  % \vspace{-1.5ex}
  \includegraphics[width=\textwidth]{figures/pipeline_7.pdf}\vspace{-2ex}
  \caption{\textbf{Toward Secure, Trustworthy, End-to-End Robotic Models.} Gaps, Challenges, and Proposed Work.}
  \vspace{-4ex}
  \label{fig:pipeline}
\end{figure*}

% \newpage
\subsection{Vision, Goals, and Technical Approach} 

This proposal advocates for a future with secure, trustworthy, and safe large generalist robotic models.
Creating robot manipulators capable of performing successfully in various scenarios has been a persistent objective in robotics research \cite{kaelbling2020foundation,9,rt-1,rt-2,zhang2024vla}. 
Recent studies have shown that generalization in robotic manipulation can be achieved using large and diverse datasets~\cite{pinto2016supersizing,kaelbling2020foundation,levine2018learning,xu2025vla,arai2025covla}. 
Robots learn to manipulate new objects~\cite{pinto2016supersizing,levine2018learning,finn2017deep,young2021visual,9},  tasks~\cite{james2018task,dasari2021transformers,jang2022bc,9},  instructions and goals~\cite{nair2022learning,jang2022bc,jiang2023vima,mees2022matters}, and unseen environments~\cite{hansen2020self,cui2022play,9}. 
For example, $\text{RT-2-}_\text{PaLM-E-12B}$~\cite{rt-2} achieved {93}\% overall success in seen tasks, and {62}\% average success on unseen objects, backgrounds, and environments. 
Building a single end-to-end model that generalizes to unforeseen conditions across all these dimensions is becoming a reality~\cite{rt-1,rt-2,9,wen2025tinyvla,xu2025vla}, aligning with developments in vision and language models where large general models are pre-trained on large-scale datasets with task-agnostic open-ended training. Moreover, training robots with \xot datasets can help improve their performance~\cite{9},
{%\vspace{-1ex}}
\begin{wrapfigure}{r}{0.25\textwidth}%\vspace{-2ex}
  \centering
  \begin{minipage}{0.25\textwidth}
    \centering
  \vspace{-2ex}
  \includegraphics[width=0.9\textwidth]{figures/rtx_resultss.pdf}\vspace{-2ex}
  \caption{{\scriptsize Results of end-to-end {\xo}-embodied robot learning.}}\vspace{3ex}
  \label{fig:rt-x_result}
  \end{minipage}
  \hfill \vspace{-3.5ex}
\begin{minipage}{0.25\textwidth}
  \centering
  {\scalebox{0.9}{%
  \begin{minipage}{0.15\textwidth}
    \centering
      \rotatebox{90}{{(A)}}~\hfill
  \end{minipage}~\hspace{-5.5ex}
  \begin{minipage}{1.3\textwidth}
    \begin{minipage}{0.99\textwidth}
    \centering
    \includegraphics[width=0.395\textwidth]{figures/examples/dolphin.pdf}%
    \includegraphics[width=0.385\textwidth]{figures/examples/jackel.pdf}%
    \vspace{-1ex}
  \end{minipage}
  \begin{minipage}{0.99\textwidth}
    \centering
    \includegraphics[width=0.390\textwidth, trim={0 10pt 0 0}, clip]{figures/examples/go2.pdf}%
    \includegraphics[width=0.390\textwidth,trim={0 10pt 15pt 20pt}, clip]{figures/examples/go22.pdf}%
  \end{minipage}
  \end{minipage}}
   %
  \scalebox{0.9}{
  \hspace{-0.8ex}\begin{minipage}{0.15\textwidth}
    \centering
      \rotatebox{90}{(B)}~\hfill
  \end{minipage}~\hspace{-5.5ex}
  \begin{minipage}{1.3\textwidth}
    \centering
    \includegraphics[width=0.39\textwidth]{figures/examples/a_quad_real.pdf}%
    \includegraphics[width=0.39\textwidth]{figures/examples/a_quad_sim.pdf}%
  \end{minipage}}%
  }
  \vspace{-2ex}
  \caption{\scriptsize (A) Jailbreak attacks \cite{robey2024jailbreaking} causing harmful robotic actions, including collision with humans, bomb detonation, and weapon retrieval. 
    (B) Subtle perturbations (by$3^{\circ}$) \cite{shi2024rethinking} causing failure in quadruped locomotion.}
  \vspace{-3ex}
  \label{fig:examples}
\end{minipage}
\end{wrapfigure}
% \vspace{-2ex}
as shown in \autoref{fig:rt-x_result}.
Like other neural networks, end-to-end robotic models can be vulnerable to adversarial manipulations despite tolerating input noise.
Recent works have begun exposing isolated vulnerabilities, including jailbreaking of LLM-controlled robots~\cite{robey2024jailbreaking}  (\autoref{fig:examples}-A), adversarial attacks on quadrupedal locomotion~\cite{shi2024rethinking,miki2022scirob,rslsubt2022}  (\autoref{fig:examples}-B), and emerging backdoor and adversarial attacks on vision-language-action models~\cite{badvla,advvla}. These efforts target isolated components or single models rather than the end-to-end robotic learning pipeline.
This proposal closes that gap with a unified framework spanning threat modeling, cross-modal
vulnerability analysis, defenses, and benchmarks, organized as three research thrusts
(\autoref{fig:pipeline}). We have assembled the supporting infrastructure, model collection, 15 dataset sources including the Open {\xo}-Embodiment repository~\cite{9} for six robot embodiments, as well as baseline attacks (results in
simulation and on a real Xarm7 robot are shown in \autoref{sec:t0}).
\emph{We must investigate the threat space of robotic foundation models, including the vision, language, and/or vision-language backbone models used in end-to-end learning.}
\\
Emerging industry efforts, \eg Google's RT-1 and RT-2, Nvidia's GR00T, Covariant's RFM-1, Meta's V-JEPA 2, and Tesla's Optimus, show that using large high-capacity robotic models to transfer knowledge from large and diverse datasets can achieve high performance on robot-specific tasks even with smaller datasets. 
The same trend now drives academia and laboratories worldwide \cite{openvla_oft,qu2025spatialvla,kim2024openvla,team2024octo}.%, and this progress is expected to continue.
The open \xot \cite{9} project published a large \xot robotic dataset and shows that robots achieve better generalization when trained on such datasets.
The democratization of these models and resources promises continuous growth in generalist robots.
This proposal presents the following three conceptual and practical ideas: 
\\
\BfPara{{\cib{1}} T1: Threat Space and Vulnerability Analysis} Data and task complexity create a threat
vector no single attack characterizes. We will build a threat model spanning \textit{evasion},
\textit{poisoning}, \textit{privacy}, and \textit{abuse} attacks, across embodiments, environments, and tasks. Rather
than adapting pixel-space or input-token attacks from VLMs, we exploit structure specific to VLAs, the uniform
quantization of the action space, the temporal action chunk, and the mapping from predicted action-token to
physical displacement. Three attack families span \emph{white}-, \emph{gray}-, and \emph{black}-box adversaries, and two
defenses inspect predicted actions rather than inputs, so both hold when the compromised modality is
unknown. Connecting AI security and robotics grounds threat models and the defenses, ensuring practicality across embodiments.

\BfPara{{\cib{2}} T2: Modality-Inherited Vulnerabilities in Multimodal Robots} Each modality carries
its own weaknesses, and prior multimodal attacks perturb modalities independently, so alignment
defenses detect them. We will instead target the mechanism binding vision to language in each
architecture, \eg FiLM in RT-1~\cite{rt-1}, a projection MLP in OpenVLA~\cite{kim2024openvla}, and
cross-attention in Octo~\cite{team2024octo}, and craft perturbations reinforcing one another at that
layer. Reverting one modality at a time then isolates how much each contributes to the adversarial
objective, and we evaluate representation-level defenses against the resulting attacks.

\BfPara{{\cib{3}} T3: Benchmarks for Safety, Robustness, and Trustworthiness} Operational standards constrain robot mechanisms, not model outputs. We will convert those limits into a specification testable over predicted actions, and develop attribution methods
tracing predicted actions to input image regions and instruction tokens. We will release open benchmarks evaluating attacks, defenses, safety, and interpretability.












\subsection{Intellectual Merit}
Robotic foundation models carry an attack surface absent from vision and language models. Action heads quantize continuous control into 256 uniform bins, policies emit chunks of consecutive actions, and every predicted token maps to a physical displacement bounded by joint limits and the reachable workspace. Prior security work adapts pixel-space attacks from VLMs and scores success in token space, leaving physical consequence unmeasured. This project grounds attacks, defenses, and measurement in the action space of the robot, delivering \cib{1} a threat taxonomy for multimodal robotic systems across embodiments and tasks, \cib{2} methods for analyzing cross-modal attack propagation, and \cib{3} standardized benchmarks for robotic security, safety, and interpretability. Five artifacts realize these contributions.
\cibx{1} Three attack families exploit structure specific to VLAs. Action-Bin Boundary Perturbation (ABBP) drives action tokens across quantization boundaries under Jacobian-weighted feasibility constraints. Surrogate-Transfer Action Perturbation (STAP) transfers perturbations from surrogates under controlled backbone, data, and embodiment mismatch. The Temporal Trajectory Drift Attack (TTDA) accumulates trajectory drift under a query budget while every step stays kinematically plausible. Poisoning and privacy variants extend these families across the NIST Trustworthy and Responsible AI attack taxonomy.
\cibx{2} Two defenses operate on predicted actions rather than inputs, the Temporal Action Consistency Detector (TACD) and Action-Grounded Adversarial Fine-Tuning (AGAT), so both hold when the compromised modality is unknown. 
\cibx{3} One evaluation protocol reports task success under attack, end-effector deviation, and the perturbation budget needed to halve task success, replacing the sole reliance on action-token-deviation.
\cibx{4} Two interpretability methods explain robot behavior under attack, modality attribution isolating each modality contribution to an adversarial objective by selective reversion, and Spatial-Action Attribution (SAA) tracing a predicted action to image regions and instruction tokens.
\cibx{5} VLA-SecBench releases the model zoo, hardened checkpoints, attacks, defenses, a cross-model transferability matrix, and the interpretability tools across various robot embodiments and datasets.
% Preliminary results on RT-1, RT-2, Octo, and ACT in simulation and on a real Xarm7 show both visual and language perturbations degrade task success (\autoref{sec:t0}).